\documentclass[11pt]{article}
\usepackage[margin=2cm]{geometry}
\usepackage{xcolor}
\usepackage{amsmath,amssymb}
\usepackage{graphicx}

\definecolor{natexorange}{RGB}{255, 107, 53}

\title{\Huge \bfseries \color{natexorange} Informe Técnico: Arquitectura y Rendimiento}
\author{\textbf{Nombre del Autor} \\ Departamento / Institución}
\date{\today}

\begin{document}

\maketitle

\begin{abstract}
Este informe técnico describe los principios de arquitectura, eficiencia y optimización de recursos para entornos de trabajo científicos y académicos.
\end{abstract}

\section{Resumen Ejecutivo}
El presente documento demuestra la viabilidad y eficiencia del entorno de desarrollo local. Mediante la integración de herramientas ligeras y portables, se eliminan las dependencias complejas y los tiempos de latencia asociados a servicios externos.

\section{Marco Teórico y Modelado}
La evolución del sistema y su gasto energético se modela a través de la relación de potencia:
\begin{equation}
  P(t) = V \cdot I(t) + \int_{0}^{t} \kappa \, \Delta T(\tau) \, d\tau
\end{equation}

\section{Conclusiones y Recomendaciones}
\textbf{NaTex} proporciona una alternativa rápida, moderna y 100\% libre de cuotas o suscripciones para la redacción de informes técnicos.

\end{document}
